\section{The Capacity-Radius Regime}

\subsection{Normalized target}

For rate \(\rho=k/n\), the list-decoding capacity benchmark is
\[
  1-\rho.
\]
At finite block length we use an explicitly strict integer radius
\[
  r_{\mathrm{cap}}<n-k,
  \qquad
  \deltacap=\frac{r_{\mathrm{cap}}}{n}<1-\rho.
\]
The strict inequality separates the intended capacity regime from endpoint
rounding questions.

For the running rate-\(1/4\), length-\(4096\) example,
\[
  n=4096,\qquad k=1024,\qquad r_{\mathrm{cap}}=3071,
\]
and hence
\[
  \deltacap=\frac{3071}{4096}<\frac34.
\]
This example is a finite parameter checkpoint, not the sufficiently-large-field
family required for a fixed-bit security claim.

\subsection{Field-size scaling}

Fix a target error fraction \(\epsstar\), such as \(2^{-128}\).  A quantitative
list or MCA theorem must produce a bound \(B\) satisfying
\[
  B\le\epsstar|\F|.
\]
For an integer-valued nonempty list, this already forces
\(|\F|\ge1/\epsstar\).  Therefore a small fixed field cannot instantiate the
grand list-decoding target at cryptographic \(\epsstar\), regardless of metric
normalization.

The main theorem must consequently be parameterized by a family
\[
  (\F_\lambda,\mathcal L_\lambda,n_\lambda,k_\lambda)
\]
with constant rate, smooth evaluation domains, and field size growing
sufficiently quickly relative to the derived list-size or bad-challenge bound.

\subsection{Largest safe radius}

For fixed \(\mathcal C,m,\epsstar\), define the list-decoding safe-radius set
\[
  \mathcal R_{\mathrm{list}}
  =
  \left\{
    \delta\in[0,1]:
    L_m(\mathcal C,\delta)\le\epsstar|\F|
  \right\}.
\]
A grand-challenge solution must not merely exhibit one member of this set; it
must determine or tightly characterize its largest element.  The first proof
milestone in this manuscript is more modest: establish that a strict
capacity-approaching radius belongs to this set under explicit field-size
hypotheses.
